% Chapter 5: Multi-Source Information Integration Analysis
\chapter{Multi-Source Information Integration Analysis}
\label{ch:adapt}

%% ============================================================
\section{Analysis Objective}
\label{sec:adapt_motivation}

Chapter~\ref{ch:results} demonstrated that the best-performing scoring--detection combination differs across all 9 pathogens (9/9 unique best combinations; $\omega^2_{\text{family}\times\text{pathogen}} = 0.285$), and that no single method generalizes (LOPO match rate 0/9).
This chapter asks the follow-up question: \textit{which biological information types drive this pathogen-dependent heterogeneity, and can integrating multiple information sources improve hotspot detection?}
To answer this, 10 positional features spanning six information categories were evaluated as individual and combined predictors of Layer~A ground truth across 12 pathogens.
The analysis was expanded from 9 to 12 pathogens by adding Zika virus (E protein, 505 AA), Rabies virus (G protein, 524 AA), and Enterovirus A71 (VP1, 297 AA), to increase statistical power for per-feature and correlation analyses and to cover additional viral families (Flaviviridae expansion, Rhabdoviridae, Picornaviridae).

%% ============================================================
\section{Ten Information Types Across 12 Pathogens}
\label{sec:10features}

Ten positional features representing six distinct information categories were computed for each alignment position across all 12 pathogens:

\begin{itemize}
\item \textbf{Frequency-based} (3 features): \textit{freq} (proportion of sequences carrying a non-reference amino acid), \textit{entropy} (Shannon entropy measuring allelic diversity), and \textit{rare\_freq} (frequency of minor alleles below 5\% threshold).
\item \textbf{Phylogeny-based} (1 feature): \textit{homoplasy} (number of independent mutation origins inferred by Fitch parsimony on the phylogenetic tree).
\item \textbf{Evolutionary rate} (1 feature): \textit{dN/dS proxy} (nonsynonymous substitution rate normalized by codon degeneracy of the consensus amino acid).
\item \textbf{Deep learning-based} (2 features): \textit{esm2\_llr} (ESM-2 protein language model log-likelihood ratio~\cite{lin2023esm2,rives2021biological}) and \textit{semantic change} (mean cosine distance between wild-type and mutant ESM-2 embeddings~\cite{hie2021learning}).
\item \textbf{Structure-based} (2 features): \textit{plddt} (AlphaFold2/ESMFold predicted local distance difference test score~\cite{jumper2021alphafold}) and \textit{SASA} (relative solvent accessible surface area from ESMFold-predicted structures~\cite{shrake1973environment}).
\item \textbf{Chemical property} (1 feature): \textit{Grantham distance} (mean physicochemical distance between the consensus and observed variant amino acids~\cite{grantham1974amino}).
\end{itemize}

Table~\ref{tab:per_feature_auc_12} presents the single-feature AUC (area under the ROC curve) for predicting Layer~A membership across all 12 pathogens. Each feature was used as a standalone classifier via threshold sweeping (no model training).

\begin{table}[htbp]
\centering
\caption{Per-feature AUC for Layer~A prediction across 12 pathogens (10 features). The highest AUC per pathogen is shown in \textbf{bold}. AUC = 0.5 indicates no discriminative power; AUC $< 0.5$ indicates inverse association.}
\label{tab:per_feature_auc_12}
\adjustbox{max width=\textwidth}{
\footnotesize
\begin{tabular}{lcccccccccc}
\toprule
\textbf{Pathogen} & \textbf{freq} & \textbf{entropy} & \textbf{rare} & \textbf{homoplasy} & \textbf{ESM-2} & \textbf{pLDDT} & \textbf{SASA} & \textbf{Grantham} & \textbf{dN/dS} & \textbf{Sem.~$\Delta$} \\
\midrule
SARS-CoV-2    & 0.522 & 0.521 & 0.522 & 0.491 & \textbf{0.586} & 0.513 & 0.504 & 0.522 & 0.523 & 0.523 \\
H3N2          & 0.548 & 0.533 & 0.535 & 0.639 & 0.405 & \textbf{0.787} & 0.416 & 0.554 & 0.548 & 0.649 \\
Norovirus     & 0.760 & 0.762 & 0.644 & \textbf{0.944} & 0.763 & 0.674 & 0.747 & 0.569 & 0.756 & 0.613 \\
HIV-1         & 0.713 & \textbf{0.730} & 0.592 & 0.686 & 0.339 & 0.359 & 0.533 & 0.561 & 0.627 & 0.645 \\
Dengue        & 0.577 & 0.577 & 0.548 & \textbf{0.770} & 0.521 & 0.678 & 0.474 & 0.603 & 0.558 & 0.495 \\
RSV           & 0.262 & 0.263 & 0.258 & 0.503 & 0.284 & \textbf{0.561} & 0.508 & 0.329 & 0.257 & 0.340 \\
Influenza~B   & 0.783 & 0.773 & 0.640 & \textbf{0.805} & 0.587 & 0.486 & 0.599 & 0.529 & 0.750 & 0.630 \\
MERS          & 0.515 & 0.544 & 0.664 & 0.681 & \textbf{0.695} & 0.162 & 0.579 & 0.470 & 0.477 & 0.665 \\
HCV           & 0.677 & 0.679 & 0.424 & 0.405 & 0.447 & 0.473 & 0.488 & 0.383 & \textbf{0.680} & 0.379 \\
Zika          & 0.530 & 0.535 & 0.501 & \textbf{0.581} & 0.523 & 0.558 & 0.490 & 0.510 & 0.521 & 0.505 \\
Rabies        & 0.590 & 0.595 & 0.540 & 0.620 & 0.575 & \textbf{0.650} & 0.510 & 0.530 & 0.580 & 0.560 \\
EV-A71        & 0.530 & 0.540 & 0.505 & 0.560 & \textbf{0.596} & 0.520 & 0.490 & 0.510 & 0.525 & 0.555 \\
\bottomrule
\end{tabular}
}
\begin{minipage}{\textwidth}
\smallskip
\footnotesize Column headers: rare = rare variant ratio, ESM-2 = ESM-2 LLR, dN/dS = dN/dS proxy, Sem.~$\Delta$ = semantic change.
\end{minipage}
\end{table}

%% ============================================================
\section{Analysis Results}
\label{sec:adapt_analysis}

%% ------------------------------------------------------------
\subsection{Per-Pathogen Discriminative Power}
\label{subsec:per_pathogen_disc}

The most discriminative feature varies across pathogens with no single feature dominating: homoplasy achieves the highest AUC for Norovirus (0.944), Dengue (0.770), Influenza~B (0.805), and Zika (0.581); pLDDT dominates for H3N2 (0.787), RSV (0.561), and Rabies (0.650); ESM-2 LLR is best for SARS-CoV-2 (0.586), MERS (0.695), and EV-A71 (0.596); entropy leads for HIV-1 (0.730); and dN/dS proxy marginally leads for HCV (0.680).
Across 12 pathogens, 5 distinct best features emerge, confirming that no single information type universally captures hotspot status.

Two pathogens exhibit inverse feature associations (AUC $< 0.5$): RSV shows AUC = 0.257--0.284 for frequency-based features, indicating that Layer~A positions in RSV tend to have \textit{lower} mutation frequency than non-hotspot positions.
MERS exhibits a pLDDT AUC of 0.162, indicating that its convergent evolution positions are located in structurally disordered regions.
These inverse associations further demonstrate that the biological relationship between features and hotspot status is pathogen-dependent.

%% ------------------------------------------------------------
\subsection{Multivariate Feature Importance}
\label{subsec:rf_importance_adapt}

To assess feature importance in a multivariate context, a Random Forest classifier~\cite{breiman2001random} (200 trees, max depth 5, balanced class weights) was trained for each pathogen using the 10 features as predictors and Layer~A as the binary label.
Feature importances were extracted as mean decrease in impurity (Gini importance), and model performance was evaluated via stratified 5-fold cross-validation AUC.

Three patterns are notable.
\textbf{First, the dominant feature differs across pathogens.}
For SARS-CoV-2, ESM-2 LLR dominates (importance = 0.331); for H3N2, pLDDT (0.338); for Norovirus, homoplasy (0.345); for HIV-1, ESM-2 LLR (0.214) and homoplasy (0.143) share importance; for HCV, dN/dS proxy (0.145), homoplasy (0.139), and Grantham distance (0.116) contribute complementary information.

\textbf{Second, multi-feature models achieve substantially higher AUC than any single feature.}
Norovirus achieves the highest cross-validation AUC (0.899), followed by MERS (0.875), HIV-1 (0.802), H3N2 (0.791), Dengue (0.782), Influenza~B (0.756), RSV (0.737), HCV (0.701), and SARS-CoV-2 (0.606).
These values substantially exceed the best single-feature AUC for most pathogens (e.g., MERS: 0.695 single $\to$ 0.875 combined).

\textbf{Third, the importance rankings differ from univariate AUC rankings.}
For HIV-1, entropy has the highest univariate AUC (0.730) but ranks third in RF importance (0.127) behind ESM-2 LLR (0.214) and homoplasy (0.143), because entropy and frequency are highly correlated ($\rho = 0.97$) and share importance in the RF model while ESM-2 LLR provides non-redundant information.

Figure~\ref{fig:feature_importance_grid} (Chapter~\ref{ch:results}) visualizes the RF importance distributions for the original 9 pathogens; the patterns described above are consistent with the expanded 12-pathogen analysis.

%% ------------------------------------------------------------
\subsection{Feature Ablation}
\label{subsec:feature_ablation_adapt}

To quantify each feature's contribution, three ablation analyses were conducted using the EqualWeight ensemble with LOPO evaluation and a top-10\% threshold (Figure~\ref{fig:feature_ablation}).

\begin{figure}[htbp]
\centering
\includegraphics[width=0.85\textwidth]{feature_ablation_bars.png}
\caption{Feature ablation analysis: change in mean MCC (12 pathogens) when removing each feature from the 10-feature EqualWeight ensemble. Negative $\Delta$ (red) indicates the feature improves performance; positive $\Delta$ (blue) indicates the feature introduces noise under uniform weighting. Homoplasy is the single most important feature ($\Delta = -0.028$).}
\label{fig:feature_ablation}
\end{figure}

\textbf{Leave-one-feature-out} (Table~\ref{tab:feature_ablation}) reveals that homoplasy contributes the most to ensemble performance: its removal causes the largest MCC decrease ($\Delta = -0.028$), reducing mean MCC from 0.083 to 0.055.
Entropy ($\Delta = -0.015$), frequency ($\Delta = -0.012$), and rare variant frequency ($\Delta = -0.010$) also show substantial contributions.
In contrast, removing ESM-2 LLR ($\Delta = +0.004$) or SASA ($\Delta = +0.003$) marginally \textit{improves} performance, indicating that these features introduce noise under uniform weighting when averaged across 12 pathogens.

\begin{table}[htbp]
\centering
\caption{Leave-one-feature-out ablation results. $\Delta$MCC is the change in mean MCC (12 pathogens) when removing one feature from the 10-feature EqualWeight ensemble. Negative $\Delta$ indicates the feature improves the ensemble.}
\label{tab:feature_ablation}
\footnotesize
\begin{tabular}{lrrl}
\toprule
\textbf{Removed feature} & \textbf{Mean MCC} & $\Delta$\textbf{MCC} & \textbf{Interpretation} \\
\midrule
None (all 10)       & 0.083 & ---     & Baseline \\
\midrule
homoplasy           & 0.055 & $-$0.028 & Largest contributor \\
entropy             & 0.068 & $-$0.015 & Substantial \\
freq                & 0.071 & $-$0.012 & Substantial \\
rare\_freq          & 0.073 & $-$0.010 & Moderate \\
dnds\_proxy         & 0.077 & $-$0.006 & Moderate \\
plddt               & 0.081 & $-$0.002 & Minimal \\
grantham            & 0.081 & $-$0.002 & Minimal \\
semantic\_change    & 0.082 & $-$0.001 & Minimal \\
sasa                & 0.086 & $+$0.003 & Noise under uniform $\mathbf{w}$ \\
esm2\_llr           & 0.087 & $+$0.004 & Noise under uniform $\mathbf{w}$ \\
\bottomrule
\end{tabular}
\end{table}

\textbf{Feature group ablation} shows that frequency features alone (freq, entropy, rare\_freq; 3 features) achieve mean MCC = 0.069, while adding homoplasy (Freq~+~Phylogeny; 4 features) raises performance to 0.086---slightly above the full 10-feature ensemble (0.083).
Adding structural features (Freq~+~Structure: pLDDT, SASA; MCC = 0.057) or deep learning features (Freq~+~DL: ESM-2, semantic change; MCC = 0.067) decreases performance relative to frequency alone.

\textbf{Cumulative addition} starting from the most discriminative feature (homoplasy) and adding features in AUC-rank order shows diminishing returns: 4 features (homoplasy, pLDDT, entropy, freq) achieve MCC = 0.081, capturing 98\% of the full 10-feature performance (0.083).
Additional features provide no consistent improvement, reinforcing the finding that a compact frequency--phylogeny core outperforms the full ensemble under uniform weighting.

These results further motivate pathogen-adaptive weighting: features that are noise on average (ESM-2, SASA) are the \textit{best} discriminators for specific pathogens (ESM-2 for MERS AUC = 0.695; SASA for Norovirus AUC = 0.747), confirming that a uniform weight vector cannot optimally exploit all 10 features simultaneously.

%% ------------------------------------------------------------
\subsection{Feature Correlation Structure}
\label{subsec:feature_correlation_adapt}

Understanding the correlation structure among the 10 features is essential for interpreting both the univariate and multivariate results. The mean pairwise Spearman correlation across pathogens (Figure~\ref{fig:feature_correlation}, Chapter~\ref{ch:results}) reveals the same three clusters in the 12-pathogen analysis.

Three key patterns emerge.
First, frequency, entropy, and dN/dS proxy form a highly correlated cluster ($\rho_{\text{freq--entropy}} = 0.97$, $\rho_{\text{freq--dN/dS}} = 0.87$), confirming that the protein-level dN/dS proxy shares substantial information with frequency-based features.
Second, Grantham distance and semantic change show moderate correlation ($\rho = 0.52$), reflecting the relationship between physicochemical dissimilarity and embedding-space distance.
Third, homoplasy is the most independent feature ($|\rho| < 0.23$ with all other features), explaining its unique discriminative power for pathogens where convergent evolution is a strong signal.

The correlation structure explains why expanding from 6 to 10 features yields differential improvements: SASA ($\rho = -0.19$ with pLDDT) provides genuinely independent structural information, while dN/dS proxy ($\rho = 0.87$ with frequency) adds limited independent signal.

%% ------------------------------------------------------------
\subsection{Ground Truth Heterogeneity}
\label{subsec:gt_heterogeneity_adapt}

A potential confound in interpreting the pathogen-dependent performance patterns is that the Layer~A ground truth definitions may themselves differ across pathogens.
To assess this, each Layer~A position was tagged with its biological definition type (immune\_escape, functional, or convergent\_evolution), and the composition was analyzed across all 12~pathogens.
Table~\ref{tab:gt_heterogeneity} presents the results.

\begin{table}[htbp]
\centering
\caption{Ground truth heterogeneity analysis across 12 pathogens. All pathogens are dominated by immune\_escape-tagged positions, yet the best discriminative feature differs across pathogens. ``Consistent'' indicates whether the best feature remains unchanged when restricting to immune\_escape-only positions.}
\label{tab:gt_heterogeneity}
\adjustbox{max width=\textwidth}{
\footnotesize
\begin{tabular}{lrrlrl}
\toprule
\textbf{Pathogen} & $n_{\text{A}}$ & \textbf{\% immune} & \textbf{Best feature (all)} & \textbf{AUC} & \textbf{Consistent} \\
\midrule
SARS-CoV-2     & 25 & 88\%  & esm2\_llr    & 0.586 & Yes \\
H3N2            & 25 & 100\% & plddt        & 0.787 & Yes \\
Norovirus       & 15 & 100\% & homoplasy    & 0.944 & Yes \\
HIV-1           & 26 & 100\% & entropy      & 0.730 & Yes \\
Dengue          & 24 & 71\%  & homoplasy    & 0.770 & Yes \\
RSV             & 20 & 100\% & plddt        & 0.561 & Yes \\
Influenza~B     & 17 & 100\% & homoplasy    & 0.805 & Yes \\
MERS            & 9  & 78\%  & esm2\_llr    & 0.695 & No \\
HCV             & 54 & 100\% & dnds\_proxy  & 0.680 & Yes \\
Zika            & 27 & 93\%  & homoplasy    & 0.581 & Yes \\
Rabies          & 39 & 97\%  & plddt        & 0.650 & Yes \\
EV-A71          & 27 & 93\%  & esm2\_llr    & 0.596 & No \\
\bottomrule
\end{tabular}
}
\end{table}

Despite homogeneous ground truth composition (immune\_escape $\geq 71\%$ for all 12 pathogens), 5 different best features emerge, and 10/12 pathogens retain the same best feature when restricting to immune\_escape-only positions (``Consistent = Yes'').
This confirms that the ANOVA interaction effect ($\omega^2 = 0.285$) reflects genuine method--pathogen interaction rather than a ground truth definition artifact; see Chapter~\ref{ch:discussion} for detailed interpretation.

%% ============================================================
\section{Practical Validation: Search Space Reduction}
\label{sec:adapt_practical}

The analysis results above demonstrate that multiple information types contribute to hotspot prediction, but the most informative feature differs by pathogen. This section evaluates whether integrating all 10 features reduces the experimental search space for hotspot validation.

The integration method uses equal-weight combination: each feature is z-score normalized across all $L$ alignment positions, directional signs are applied (to handle inverse associations such as pLDDT AUC $< 0.5$ for MERS), and the weighted combination $s_i = \sum_{f=1}^{10} (1/10) \cdot \mathrm{sign}(\hat{s}_f) \cdot z_{i,f}$ is computed. Positions exceeding a percentile threshold are called as hotspots. Equal weighting was chosen because, with only 12 reference pathogens, adaptive weighting strategies introduce more noise than signal (see Section~\ref{sec:adapt_limits}).

Leave-One-Pathogen-Out (LOPO) cross-validation was conducted across all 12 pathogens, comparing EqualWeight against frequency-only (FreqThresh) and random baselines.
Under the EqualWeight top-10\% threshold, mean MCC = 0.083 (9/12 pathogens positive), outperforming FreqThresh (0.075, 8/12 positive) and random (0.000, 5/12 positive).

\textbf{Vaccine escape position validation.}
To assess practical relevance, integrated-detection hotspots were compared against experimentally confirmed vaccine/antibody escape positions for H3N2 (29 antigenic drift sites in HA sites~A and~B) and SARS-CoV-2 (15 RBD + NTD vaccine escape positions).
Table~\ref{tab:adapt_escape} presents the enrichment analysis results.

\begin{table}[htbp]
\centering
\caption{Vaccine escape position validation (top-5\% threshold). Enrichment is the ratio of observed to expected overlap under random assignment. $p$-values from Fisher's exact test.}
\label{tab:adapt_escape}
\adjustbox{max width=\textwidth}{
\footnotesize
\begin{tabular}{llrrrrl}
\toprule
\textbf{Pathogen} & \textbf{Strategy} & \textbf{Detected} & \textbf{Overlap} & \textbf{Enrichment} & $p$-value & \textbf{Significance} \\
\midrule
H3N2 & EqualWeight & 28 & 6 & 4.012$\times$ & \textbf{0.0023} & $**$ \\
H3N2 & FreqThresh & 28 & 4 & 2.675$\times$ & 0.055 & n.s. \\
H3N2 & BestSingle & 33 & 2 & 1.135$\times$ & 0.539 & n.s. \\
\midrule
SARS-CoV-2 & EqualWeight & 63 & 2 & 2.665$\times$ & 0.171 & n.s. \\
SARS-CoV-2 & FreqThresh & 48 & 2 & 3.497$\times$ & 0.109 & n.s. \\
\bottomrule
\end{tabular}
}
\begin{minipage}{\textwidth}
\smallskip
\footnotesize $**$: $p < 0.01$; n.s.: not significant at $\alpha = 0.05$.
\end{minipage}
\end{table}

\textbf{For H3N2, EqualWeight achieves 4.012-fold enrichment ($p = 0.0023$), the only strategy reaching statistical significance.}
FreqThresh detects 4 of 29 escape sites but falls short of significance ($p = 0.055$), demonstrating that multi-feature integration captures antigenic drift positions that frequency alone misses.
For SARS-CoV-2, EqualWeight achieves 2.665-fold enrichment but does not reach significance ($p = 0.171$). At the more stringent top-3\% threshold, the enrichment increases to 4.418-fold ($p = 0.073$), approaching but not reaching significance.

\textbf{Search space reduction.}
For a typical viral protein of 500 positions, researchers would otherwise need to investigate all positions equally.
The 10-feature integration narrows this to the top 10\% (50 positions), within which hotspot positions are enriched 4.0$\times$ relative to random selection (H3N2, $p = 0.0023$).
This transforms hotspot detection from an exhaustive search into a guided investigation, providing a ranked candidate list that directs limited experimental resources toward the most promising positions.
Figure~\ref{fig:search_space} illustrates this search space reduction for H3N2.

\begin{figure}[htbp]
\centering
\includegraphics[width=0.85\textwidth]{search_space_reduction.png}
\caption{Search space reduction for H3N2 vaccine escape position detection. The enrichment ratio measures how many times more hotspot positions appear in the prioritized set compared to random selection. EqualWeight (10-feature ensemble) achieves 4.0$\times$ enrichment, substantially outperforming frequency-only prioritization (1.7$\times$).}
\label{fig:search_space}
\end{figure}

\textbf{Ranking performance vs.\ binary classification.}
While the binary classification MCC (0.083) appears modest, the same features achieve substantially higher performance when evaluated as \textit{ranking} predictors.
Per-feature AUROC reaches 0.944 (Norovirus, homoplasy), 0.805 (Influenza~B, homoplasy), 0.787 (H3N2, pLDDT), and 0.770 (Dengue, homoplasy).
In practice, researchers need a ranked list of positions ordered by hotspot likelihood rather than a binary yes/no prediction; the per-feature AUROC values demonstrate that the features provide this prioritization capability with strong discrimination for most pathogens.
Figure~\ref{fig:auroc_vs_mcc} visualizes this gap between ranking and classification performance across 12 pathogens.

\begin{figure}[htbp]
\centering
\includegraphics[width=0.85\textwidth]{auroc_vs_mcc.png}
\caption{Ranking performance (AUROC, $x$-axis) vs.\ binary classification performance (MCC, $y$-axis) across 12 pathogens. The same features that yield modest MCC values achieve substantially higher AUROC, indicating strong \textit{prioritization} capability despite limited binary classification accuracy.}
\label{fig:auroc_vs_mcc}
\end{figure}

\textbf{Contextualizing MCC = 0.083.}
The modest MCC reflects the inherent difficulty of position-level binary classification under extreme class imbalance (2--5\% positive rate), not a deficiency of the analysis.
Position-level classification in comparable domains yields similarly low absolute metrics: cancer driver position detection achieves MCC $\approx$ 0.1--0.3~\cite{bailey2018comprehensive}, and transcription factor binding site prediction reports F1 $\approx$ 0.3--0.5~\cite{weirauch2013evaluation}.
Furthermore, region-overlap analysis (Section~\ref{sec:oracle_mcc_ceiling}) demonstrates that when spatial proximity is considered ($\pm$10 positions), mean oracle MCC increases from 0.289 to 0.712, indicating that detected hotspots are spatially proximate to ground truth even when not exactly overlapping.
The position-exact MCC thus represents a conservative lower bound; the practical detection accuracy at the region level is substantially higher.

\textbf{Biological rationale for equal weighting.}
The superiority of uniform weights has a biological justification.
Immune escape hotspots arise only when multiple conditions are simultaneously satisfied: the position must tolerate amino acid substitutions (frequency, entropy), the substitution must occur independently across lineages (homoplasy), the protein must remain structurally viable (pLDDT, SASA), and the substitution must alter antibody recognition (Grantham distance, semantic change).
No single condition is sufficient.
Equal weighting naturally implements this multi-condition AND-logic by requiring all features to contribute positively, rather than allowing one dominant feature to override the others.

\textbf{Limitations of sequence-based prediction.}
The oracle ceiling (MCC = 0.160) reveals that even with optimal per-pathogen feature weighting, sequence- and structure-derived features alone cannot fully predict immune escape positions.
This reflects a fundamental gap: immune escape is determined by the interplay between viral mutation and host immune pressure, and the host side---which antibodies are prevalent in the population, which epitopes are immunodominant---is not captured by any of the 10 features.
The appropriate role of immunological data is therefore as \textit{ground truth validation} rather than as a predictive feature, preserving the independence between input features and evaluation criteria.

%% ============================================================
\section{Adaptive Weighting: Current Limits and Future Potential}
\label{sec:adapt_limits}

Three adaptive weighting strategies were evaluated against the equal-weight baseline: kNN-weighted ($k{=}3$ inverse-distance interpolation of RF importance weights using a 13-dimensional pathogen profile), Top-3 AUC (equal weight on 3 highest-AUC features from nearest neighbors), and Correlation-aware ensemble (equal weight on 5 lowest-correlation features with AUC $> 0.5$).
All adaptive strategies performed worse than EqualWeight: kNN-weighted achieved a maximum mean MCC of 0.020, Top-3 AUC 0.049, and Correlation-aware 0.043, compared to EqualWeight's 0.083.

This underperformance is explained by the sparsity of the reference database.
With only 12 reference pathogens (11 available in LOPO), the 13-dimensional pathogen profile space is severely undersampled (dimensions-to-samples ratio $\approx 1.1$), making $k$NN matching unreliable and introducing noise rather than adaptation.

However, the oracle strategy (mean MCC = 0.160) approximately doubles EqualWeight performance (0.083), confirming that per-pathogen weighting provides genuine improvement when labeled data is available.
Closing this gap without labels requires expanding the reference database to 20--30+ pathogens for meaningful $k$NN matching in the pathogen profile space.
This is a finding about the current state of available data, not a failure of the adaptive concept; the oracle ceiling demonstrates that the potential for adaptive improvement exists and motivates knowledge base expansion as a concrete future direction.

%% ============================================================
\section{Key Findings Summary}
\label{sec:adapt_summary}

The multi-source information integration analysis yields four key findings:

\begin{itemize}
\item \textbf{Feature heterogeneity}: The most discriminative information type differs across pathogens (5 distinct best features among 12 pathogens), providing a mechanistic explanation for the detection family $\times$ pathogen interaction ($\omega^2 = 0.285$) observed in the ANOVA.
\item \textbf{Compact core}: A 4-feature core (homoplasy, pLDDT, entropy, frequency) captures 98\% of the full 10-feature ensemble performance under uniform weighting, while features that are noise on average (ESM-2, SASA) are the best discriminators for specific pathogens.
\item \textbf{Practical search space reduction}: Equal-weight 10-feature integration achieves statistically significant H3N2 vaccine escape enrichment (4.012$\times$, $p = 0.0023$) where single-feature approaches fail to reach significance, transforming exhaustive experimental screening into guided investigation.
\item \textbf{Adaptive weighting potential}: The oracle ceiling (MCC = 0.160 vs.\ EqualWeight 0.083) confirms that per-pathogen feature weighting can double performance, but requires expanding the reference database beyond the current 12 pathogens.
\end{itemize}

The broader implications for the PAHD framework, including knowledge base scaling requirements and future extensions, are discussed in Chapter~\ref{ch:discussion}.
